% ============================================
% 第1章：机器学会了搭桥，但它不知道自己在搭桥
% 梁桥 · 最朴素的跨越
% ============================================

\chapter{机器学会了搭桥，但它不知道自己在搭桥}

\vspace{0.5cm}

\begin{center}
{\small\color{muted} 梁桥是最简单的桥。一根梁，两个支座，就跨越了。}
{\small\color{muted} 但最简单的，往往最难做好。}
\end{center}

\vspace{1cm}

\section{AI是怎么学习的？}

2026年7月2日，我在准备一场报告。报告的名字叫“AI时代，教学的变与不变”。我对着录音笔说了一段话，后来被转成了文字：

\begin{shiyu-inline}
“AI的学习方式咱们可以借鉴。它是怎么学习的呢？就是通过训练，给它大量的文本。很多时候它也不理解这东西是啥，或者说我们也很难说明什么是'理解'。在这种情况下，就不能说明什么是理解。那就是输入大量的文本，在里面去进行贝叶斯网络的转换。可能这个以后就是未来的理解，到底什么是理解呢？那么这就是理解。”

\hfill ——2026年7月2日，录音转写
\end{shiyu-inline}

这段话里有一个词，我自己说出来的时候都没意识到它有多重要：\textbf{“它也不理解这东西是啥。”}

AI的学习方式，拆开来看，其实很简单。

\subsection{第一步：喂数据}

你给AI看海量的文本。桥梁工程教材、设计规范、施工方案、论文、博客、论坛帖子\ldots\ldots 它看过的桥，比任何桥梁工程师一辈子见到的都多。

\subsection{第二步：找规律}

AI在这些文本中寻找统计规律。它发现“简支梁”这个词后面经常跟着“跨中弯矩”；它发现“$qL^2/8$”这个公式常常出现在“简支梁”和“弯矩”附近；它发现“悬索桥”和“主缆”“锚碇”“吊杆”这些词有很高的共现概率。

\subsection{第三步：生成输出}

当你问它“简支梁跨中弯矩是多少”，它做的事情不是回忆、不是推导、不是理解——它是\textbf{计算}：在它见过的所有文本里，最有可能的、最合理的、最符合统计规律的输出，就是“$qL^2/8$”。

它不知道$q$是荷载，不知道$L$是跨度，不知道那个$1/8$是从哪里来的。它只知道——这个答案，命中的概率最高。

\subsection{AI的“记忆”}

在同一段录音里，我还说了另一句话：

\begin{shiyu-inline}
“机器人有上下文，我们经常说学生记得这个就忘了那个。用AI的话说，就是上下文比较短，执行力比较差，情绪价值要求比较高，月费好几千。”

\hfill ——2026年7月2日，同一段录音
\end{shiyu-inline}

AI的“记忆”有一个硬边界，叫\textbf{上下文窗口}。你可以把它想象成一个只能装固定数量文字的盒子。盒子满了，旧的被挤出去。AI只能“看见”盒子里的东西。

这个盒子在2025年大约能装15万字，2026年能装100万字。但不管多大，它仍然是一个盒子。AI没有“经验”。它不会说“我记得大一那年，我的结构力学老师说过\ldots\ldots”它不会说“上次我做这个的时候，犯了一个错误，这次我不会再犯。”

\begin{insight}
\label{insight:chap1}

\textbf{AI的学习，本质上是统计推断。}它的目标函数只有一个：让预测误差最小化。

它不在乎“桥”是什么。它不知道塔科马大桥在1940年倒塌的时候，工程师们震惊地发现，他们引以为傲的静力分析在风的作用下完全失效了。它不知道那些失败改变了一整个学科。

AI知道事实。它不知道事实的重量。
\end{insight}

\section{人是怎么学习的？}

2025年11月4日，下午3点25分，我在课堂上讲涡振原理。

\begin{classroom-scene}
\label{scene:chap1-vortex}

{
\noindent \textbf{2025年11月4日，大连交通大学，桥梁抗风课。}

\vspace{0.3cm}

\noindent 我在电子白板上放了一个CFD模拟视频。一个圆柱体，均匀的风吹过来，尾部交替脱落着涡。我指着屏幕上周期性波动的升力曲线，说：

\noindent “这个，就是涡振。”

\noindent 然后我停了一下，说了一段话，后来转写出来是这样的：

\vspace{0.3cm}

\begin{quote}
“其实我不觉得任何这个荷载特性啊，或者是风震，是你们必须学的。我更希望就是你们看到这个研究的背后——其实就是对一个简单的动力学方程，今天怎么处理，怎么去简化。把这种复杂的问题简化到那去，包括怎么去面对这么复杂的这个问题，怎么把它分解成几个独立的振动形式，然后再去讨论\ldots\ldots 所以我觉得这个反而是很很有意义的。就是一个最简单的公式，怎么能 cover 到这么多的这个现象。”
\end{quote}

\vspace{0.3cm}

\noindent ——\href{2025/文字转写/2025-11-04\_1524\_结构工程课程之涡振（涡震）原理与特性讲解.md}{2025年11月4日，涡振原理课，录音转写}
}
\end{classroom-scene}

这段话里，我没有说“今天的教学目标是理解涡振原理”。我说的是：\textbf{怎么把一个复杂问题分解成几个独立的振动形式。}

这不是知识。这是一种\textbf{思维方式}。而思维方式，AI不会教。

\subsection{人学习时，身体在动，手心在出汗}

2026年5月22日，我在面向呼和浩特铁路局教师的线上分享会上，讲游戏化教学。我讲到了一件事，后来被转写成了文字：

\begin{shiyu-inline}
“学生在上这门课之前已经学过材料力学、结构力学、或者理论力学，但是他搭建这个桥梁的过程或者状态，跟他六岁的小孩搭建这个状态是完全一致的。我们不是说他这个多么低龄或多么幼稚。我们在说，其实我们讲授的很多的理论的这种课程，如果没有经过这种实操、试错、迭代、失败成功的这个过程，他可能很难入到你一个学生的这个直觉里面去。”

\hfill ——2026年5月22日，游戏化教学分享，录音转写
\end{shiyu-inline}

这句话里有一个特别重要的词：\textbf{直觉}。

一个学完了材料力学、结构力学、理论力学的大学生，当他第一次打开Poly Bridge搭桥的时候，他的行为和一个六岁小孩\textbf{完全一样}。他不是在调用公式，他是在用直觉——这里加一根杆试试，那里加高一点试试，塌了，换一种方式，又塌了，再换一种\ldots\ldots

这不是退化。这是\textbf{学习的本质状态}。

\begin{insight}
\label{insight:chap1-intuition}

\textbf{人的学习不是模式匹配。人的学习是——经历从不知道到知道的过程，并被这个过程改变。}

AI的输出是“答案”。人的输出是“我懂了”。AI不会在输出正确结果时心跳加速。AI不会在深夜想起自己第一次搭桥成功的那一瞬间，嘴角微微上扬。
\end{insight}

\subsection{失败的权利}

在同一场分享里，我还说了一句话：

\begin{shiyu-inline}
“学生在学习过程中从来没有失败的权利，他也就没有成功的那个喜悦或者荣耀。”

\hfill ——2026年5月22日，同一场分享
\end{shiyu-inline}

这句话，是我在讲完Poly Bridge之后说的。我给学生看了一段视频：一座桥，小车开上去，桥塌了。然后学生改了三个地方，又试了一次。桥站住了。

整个过程不到两分钟。但这两分钟里，学生经历了：\textbf{困惑→尝试→失败→观察→调整→再尝试→成功→骄傲。}

在传统的课堂里，这个回路通常被压缩成一个步骤：\textbf{“记住这个公式，考试会考。”}

\section{两种学习，两种“理解”}

\begin{table}[H]
\centering
\begin{tabular}{p{0.35\textwidth} p{0.35\textwidth}}
\toprule
\textbf{AI的学习} & \textbf{人的学习} \\
\midrule
目标：最小化预测误差 & 目标：理解世界，并能在其中行动 \\
\addlinespace
输入：海量文本 & 输入：有限但多模态的体验（视觉、触觉、情感） \\
\addlinespace
记忆：上下文窗口内的完美复制 & 记忆：选择性的，在遗忘中保留本质 \\
\addlinespace
反馈：损失函数数值下降 & 反馈：\textbf{桥塌了、老师点头了、心里“啊哈”了一声} \\
\addlinespace
错误：统计偏差，需要更多数据 & 错误：\textbf{最有价值的老师} \\
\addlinespace
“理解”：能生成正确的输出 & “理解”：能解释、能迁移、\textbf{能在自己错的时候知道自己错了} \\
\bottomrule
\end{tabular}
\caption{AI的学习 vs 人的学习}
\label{tab:ai-vs-human}
\end{table}

\section{梁桥的隐喻}

梁桥是所有桥型里最简单的。一根梁，两个支座。它不优雅，不壮观，不需要任何高级的力学原理。它只是\textbf{直接跨越}。

但梁桥有一个致命的弱点：\textbf{跨度越大，自重越大。}

2025年10月21日，我在课堂上讲到这件事：

\begin{shiyu-inline}
“我们的桥梁的载重比其实是非常小的。一般来说，一个大跨度桥梁，大概百分之八十是它的自重，百分之二十是车辆。所以说，我们实际上只有百分之二十的承载能力是用在了承载荷载上。所以这个就是很大的浪费嘛。就是姜文的电影是吧——为了这碟醋包了顿饺子。”

\hfill ——2025年10月21日，桥梁课程，录音转写
\end{shiyu-inline}

“为了这碟醋包了顿饺子”——这句话，是我课堂上说的。后来这段录音被转写成文字，我重新读到的时候，发现它比我想象的更好。

\textbf{梁桥的自重，就是AI时代教育面临的问题。}

我们用80\%的精力教学生“知识”——公式、规范、软件操作。这些东西，AI一秒钟就能生成。我们只有20\%的精力用在真正重要的东西上——\textbf{困惑、尝试、失败、骄傲、判断、审美、责任}。

就像梁桥，80\%的自重，20\%的荷载。\textbf{效率极低。}

但这不是AI的错。AI只是把这个问题暴露了出来。它把“做出来”的门槛降到了零，于是我们突然发现：\textbf{原来我们花了80\%的精力教的东西，根本不值钱。}

\begin{insight}
\label{insight:chap1-beam}

\textbf{梁桥的困境，就是当代教育的困境。}

不是AI太强了，是我们一直在教AI擅长的东西。当AI能生成一切“看起来像桥”的东西，真正稀缺的就不再是生成，而是\textbf{判断}。不再是速度，而是\textbf{责任}。不再是看起来像，而是\textbf{真的站得住}。
\end{insight}

\section{这一章，我们站到了哪里}

这本书叫《请把手弄脏》。这五个字来自一个很具体的画面：一个学生，打开电脑，自己写了一个有限元求解器，然后指着它说——“这是我的。”

这个画面出现在这本书的尾声。但在这之前，我们还有很多路要走。

这一章，我们只做了一件事：\textbf{把AI的学习和人的学习，放在一起看。}

AI的学习是统计推断。人的学习是意义生成。AI输出的是“答案”。人输出的是“我懂了”。AI不会骄傲。AI不会在深夜想起自己第一次搭桥成功的那一瞬间。

梁桥是最简单的桥。但最简单的，往往最难做好。

\begin{shiyu-note}
\label{note:chap1}

\textbf{AI已经会搭桥了，问题是人还想不想过桥。}

“AI的学习”这个词，本身就是一个陷阱。我们把“输出正确答案”和“理解”混为一谈了。就像我们把“会背乘法表”和“懂数学”混为一谈一样。

AI时代的教育，真正的问题不是“AI能替代什么”，而是\textbf{“当AI能替代一切，人还剩下什么？”}

人剩下的，恰好是AI永远无法拥有的东西：困惑、好奇心、挫败感、骄傲、责任感、审美判断、以及“非做不可”的理由。

这些，才是学习的真正内容。而AI的出现，只是帮我们把这些东西从“知识传递”的杂音中清晰地分离了出来。
\end{shiyu-note}

\vspace{1cm}

\begin{decision-point}
\label{decision:chap1}

\noindent\textbf{这一章，我们看到了两种“学习”的本质区别。}

\vspace{0.3cm}

\noindent 现在，打开你的下一节课的PPT。找到第一页。删掉“课程目标”。换成一句话，让你的学生\textbf{想}知道接下来会发生什么。

\noindent 不是“通过本章学习，你将掌握\ldots\ldots”
\noindent 而是“你有没有想过，为什么\ldots\ldots？”

\vspace{0.3cm}

\noindent 就这一句。改完发给我。
\end{decision-point}

\vspace{0.5cm}

\begin{flushright}
{\small\color{muted} 下一章：为什么孩子会蹲在沙坑里搭三个小时的桥——关于游戏的本质。}
\end{flushright}